\section{怎样在仓库里找到一段实现}

目标系统的代码都在 \filepath{source/}。测试、构建工具和教材帮助理解与验证，
但不会被编进来宾操作系统。读代码时先确定它运行在 ROM、Kernel 还是用户态，
随后再找入口和对应的公开头文件。

\begin{osgridlongtable}{L{0.18\textwidth}L{0.28\textwidth}L{0.42\textwidth}}
\textbf{目录} & \textbf{运行位置} & \textbf{先看什么}\\ \hline
\filepath{source/firmware} & 复位后的实模式/保护模式代码 & 入口汇编、链接地址、
Stage 1 读取\\ \hline
\filepath{source/kernel} & Ring 0 长模式 & \filepath{include/} 中的接口，再看
\filepath{src/} 同名模块\\ \hline
\filepath{source/user} & Ring 3 & 系统调用包装、运行时与各个 ELF 程序\\ \hline
\filepath{source/abi} & Kernel 与用户程序共同编译 & 系统调用号、结构布局和
版本常量\\ \hline
\filepath{tests} & 宿主机或 QEMU 驱动 & 哪个输入被构造，断言观察了什么\\ \hline
\filepath{tools} & 宿主机 & 镜像生成、布局审计和 QEMU 生命周期\\ \hline
\end{osgridlongtable}

\topic{头文件和源文件使用相同的模块路径}
例如页表接口位于
\filepath{source/kernel/include/os/kernel/memory/page\_table.hpp}，实现位于
\filepath{source/kernel/src/memory/page\_table.cpp}。公开契约与实现被
include/src 隔开，二者内部又都按 arch、device、fs、memory、process 等
功能分目录。替换根前缀就能从接口跳到实现。

\topic{从调用者开始，避免陷进工具函数}
想理解一次 open，先从系统调用分发找到参数和返回值，再跟到 FileTable、
FileDescription、VFS 和后端。若直接从一个字符串比较函数开始，很难知道它
服务于路径解析的哪一步。硬件路径同样如此：先找 IRQ 或请求入口，再进入端口
读写函数。

\topic{代码量只回答规模问题}
构建教材时，脚本只统计 \filepath{source/} 下的 \code{.cpp}、\code{.hpp}
和 \code{.asm}，排除测试、工具、教材和构建描述。当前结果为：

\begin{osgridlongtable}{L{0.35\textwidth}L{0.20\textwidth}L{0.33\textwidth}}
\textbf{指标} & \textbf{当前值} & \textbf{含义}\\ \hline
目标源文件 & \OsProductionFileCount & 进入固件、Kernel、ABI 或用户程序的
核心文件\\ \hline
物理行 & \OsProductionPhysicalLineCount & 包含空行和注释的原始行数\\ \hline
非空非纯注释行 & \OsProductionCodeLineCount & 固定口径下的规模切片\\ \hline
C++20 & \OsProductionCppCodeLineCount & 头文件与源文件\\ \hline
NASM Intel 汇编 & \OsProductionAssemblyCodeLineCount & 复位、模式切换和
架构入口\\ \hline
\end{osgridlongtable}

行数增加可能来自重复、失败处理或新的真实机制，单独看不出质量。三字节复位
向量对启动至关重要，数百行表格驱动代码也可能只是同一规则的展开。判断进展
仍要回到行为和测试。

\topic{遇到常量时回到协议来源}
端口号和状态位应能追到硬件资料，页表掩码应能追到 Intel 手册，盘面 offset
应能追到格式头文件与镜像工具。项目常量名说明使用位置，资料说明数值来源，
测试则检查最终字节。三者缺一时，不要只凭名字猜测。

\topic{建议的代码走读顺序}
\begin{enumerate}[label=\arabic*.]
  \item 在用户可观察行为或启动日志中选一个入口；
  \item 找到产生这条行为的最外层函数；
  \item 读对应头文件，写下参数、返回值和所有权；
  \item 沿调用链记录地址空间、锁和可能阻塞点；
  \item 对照测试中的输入与断言；
  \item 最后查看反汇编、镜像字节或 QEMU 寄存器确认硬件边界。
\end{enumerate}
